% Generated draft: requires mathematical and source review.
\documentclass[11pt]{article}
\usepackage[margin=25mm]{geometry}
\usepackage{fontspec}
\setmainfont{Latin Modern Roman}
\usepackage{amsmath,amssymb,mathtools,unicode-math}
\setmathfont{Latin Modern Math}
\usepackage{enumitem,needspace,xurl,hyperref}
\setlength{\parindent}{0pt}
\setlength{\parskip}{5pt}
\setlength{\emergencystretch}{3em}
\newcommand{\N}{\mathbb{N}}
\newcommand{\Z}{\mathbb{Z}}
\newcommand{\Q}{\mathbb{Q}}
\newcommand{\R}{\mathbb{R}}
\newcommand{\C}{\mathbb{C}}
\newcommand{\F}{\mathbb{F}}
\newcommand{\A}{\mathbb{A}}
\newcommand{\PP}{\mathbb{P}}
\newcommand{\E}{\mathbb{E}}
\newcommand{\eps}{\varepsilon}
\newcommand{\dd}{\,\mathrm{d}}
\begin{document}
\section*{UnsolvedMath ID: 30006252}
\textbf{Source problem code:} catalogue entry
\subsection{Definitions and mathematical statement}
Let \(k\) denote the base field and let \(X\) be an algebraic variety over \(k\). The supplied record does not specify \(k\), its characteristic, or whether varieties must be affine, projective, smooth, irreducible, or subject to other restrictions; these hypotheses remain unresolved. Write \(\operatorname{Aut}_k(X)\) for the group of algebraic automorphisms of \(X\) over \(k\). Define the Lie algebra of global algebraic vector fields by
\[
\operatorname{Vec}_k(X)=\Gamma\!\left(X,\operatorname{Der}_k(\mathcal O_X,\mathcal O_X)\right),
\qquad [D,E]=D\circ E-E\circ D,
\]
where the derivations are derivations of the structure sheaf over \(k\). This definition does not impose smoothness on \(X\).

The first question asks whether \(\operatorname{Aut}_k(X)\) admits a natural Lie algebra \(L_X\), together with a natural injective homomorphism of \(k\)-Lie algebras
\[
\iota_X:L_X\hookrightarrow\operatorname{Vec}_k(X).
\]
The geometric structure on the automorphism group from which \(L_X\) is to be obtained, and the precise meaning of naturality, are not specified. In particular, the notation \(\operatorname{Lie}(\operatorname{Aut}_k(X))\) requires an intended construction beyond merely viewing \(\operatorname{Aut}_k(X)\) as an abstract group. The supplied original excerpt begins with the question whether an unspecified antecedent is a subalgebra of the vector fields; identifying that antecedent with this Lie algebra follows the supplied curated statement and requires source review.

The second question asks to what extent \(X\) is determined by \(\operatorname{Aut}_k(X)\). One precise reconstruction property that could be investigated for a specified class \(\mathcal C\) of \(k\)-varieties is
\[
\forall X,Y\in\mathcal C,\qquad
\operatorname{Aut}_k(X)\cong\operatorname{Aut}_k(Y)
\ \Longrightarrow\ X\cong_k Y.
\]
Here, if the groups are interpreted as abstract groups, the isomorphism on the left is an abstract group isomorphism. A stronger reconstruction property would require every such isomorphism \(\varphi\) to be induced by an isomorphism \(f:X\to Y\), in the sense that
\[
\varphi(g)=f\circ g\circ f^{-1}
\qquad\text{for every }g\in\operatorname{Aut}_k(X).
\]
These are distinct possible meanings of reconstruction. The source does not specify the class \(\mathcal C\), the structure that group isomorphisms must preserve, or which reconstruction property is intended; neither displayed property is asserted here as a theorem or an unrestricted conjecture.

The source motivates reconstruction by reporting a theorem attributed to Whittaker (1963), labelled [Whi63]: for compact, locally Euclidean manifolds \(X,Y\), every isomorphism between their groups of continuous automorphisms is induced by a homeomorphism \(X\to Y\). It reports an analogous statement for differentiable manifolds and their diffeomorphism groups, attributed to Takens (1979) and Filipkjewich (1982), labelled [Tak79] and [Fil82], and states the corresponding corollary that automorphisms of these transformation groups are inner. These are source-reported historical statements; their complete hypotheses and bibliographic entries are not supplied. They are motivation rather than hypotheses for the algebraic questions. The excerpt explicitly cautions that comparable unrestricted algebraic reconstruction cannot be expected, citing many ``nice'' varieties with trivial automorphism group; it supplies neither examples nor a definition of ``nice'' in that algebraic remark.

The supplied dated curation labels the problem partially solved and reports that the Lie algebra embeds naturally into algebraic vector fields, that reconstruction holds in important restricted settings, and that no general reconstruction theorem exists, with trivial-automorphism examples providing an obstruction. The relevant structures, restricted settings, and exact supporting results are not identified in the record. These are supplied literature assessments, not independently verified conclusions of this draft.
\subsection{Short English statement}
The problem asks how to associate a natural Lie algebra to the algebraic automorphism group of a variety, how that Lie algebra relates to algebraic vector fields, and under what hypotheses the automorphism group determines the variety. The base field, allowed varieties, structure on the automorphism group, and strength of reconstruction remain unspecified. The supplied partially solved label and historical comparisons require mathematical and source review.
\subsection{Sources}
\begin{itemize}
\item Local UnsolvedMath registry: supplied complete source record, permanent ID \(30006252\), title \emph{Lie Algebras and Reconstruction from Variety Automorphisms}, exact source code \texttt{catalogue entry}. The record belongs to the Oberwolfach Reports Open Problems collection and supplies proposed year 2025; no individual proposer is identified.
\item Supplied source citation: \emph{Algebraic Groups} (2025), Oberwolfach Reports, pp.\ 941--1010, DOI \texttt{10.4171/owr/2025/20}. \url{https://doi.org/10.4171/owr/2025/20}. The background describes these reports as documenting mathematical workshops and their open-problem discussions. No location within the report more precise than this page range is supplied.
\item Supplied literature source, labelled ``Kraft survey and stated questions'': \url{https://doi.org/10.4171/OWR/2025/20}.
\item Supplied literature source, labelled ``Kraft--Zaidenberg, algebraically generated groups'': \url{https://doi.org/10.1112/jlms.12866}.
\item Literature assessment date: \texttt{2026-08-22}. The partially solved status and associated claims are supplied dated curation derived from the OWR record, not independently verified here. The supplied background explicitly requests independent verification of time-sensitive status and cited claims.
\item Historical references appearing in the supplied excerpt: Whittaker (1963), [Whi63]; Takens (1979), [Tak79]; and Filipkjewich (1982), [Fil82], with the last name reproduced as supplied. Full bibliographic details are absent.
\item The registry labels the statement \texttt{corrected\_verified} and reports that the official report was checked, the question isolated, and notation repaired. This is supplied source-verification metadata; this draft does not claim an independent visit to or verification of any external source.
\end{itemize}
\end{document}
